\chapter{Formal Language of Behavioural States}
\label{ch:formal-language}

\section{State, Frame, and Resistance}

\begin{definition}[Behavioural state]
A behavioural state $s \in \mathcal{S}$ is a sustained mode of action together
with its associated posture, attention pattern, local environment, and reward
expectation.
\end{definition}

\begin{definition}[Framework inertia]
The framework inertia of a state $s$ is
\[
  \FI(s) = \alpha \Reward(s) + \beta \Switch(s) + \gamma \Env(s),
\]
where $\Reward(s)$ is the immediate reward rate, $\Switch(s)$ is the cognitive
cost of loading a new task frame, $\Env(s)$ is environmental reinforcement,
and $\alpha, \beta, \gamma > 0$ weight the agent's sensitivity to each
component.
\end{definition}

\begin{remark}
This definition is useful because it separates three business-relevant levers.
If a person repeatedly fails to act, the intervention might target the reward
gradient, the cognitive startup cost, or the environment. Treating all failure
as ``motivation'' hides that design choice.
\end{remark}

\section{Target States and Barriers}

\begin{definition}[Target state]
Given a current state $s_0$, a target state $s_1$ is a behavioural state whose
expected long-run value exceeds that of $s_0$ under the agent's stated goals.
\end{definition}

\begin{definition}[Activation barrier]
The activation barrier for transition $s_0 \to s_1$ is denoted by
$\DV(s_0 \to s_1)$ and represents the minimal effort needed to interrupt the
current frame and stabilize the new one.
\end{definition}

\begin{proposition}[Barrier decomposition]
For practical planning, the barrier can be decomposed as
\[
  \DV(s_0 \to s_1) \approx b_0 + \eta_1 \FI(s_0) + \eta_2 U(s_1) - \eta_3 P,
\]
where $U(s_1)$ is startup ambiguity in the target state and $P$ is prior
preparation.
\end{proposition}

\begin{discussion}
The proposition is intentionally operational. It says that high inertia in the
present, vagueness in the future, and lack of preparation all push in the same
direction: they raise the threshold for action. The only term with the
opposite sign is preparation. That is why planning matters even before one
speaks about willpower.
\end{discussion}

\section{Worked Example: Bed to Desk}

\begin{example}[A minimal state transition]
Consider the transition from scrolling in bed to seated reading at a desk.
\begin{enumerate}[nosep]
  \item $\Reward(s_0)$ is high because scrolling offers rapid novelty.
  \item $\Switch(s_0)$ is high because the reader must remember where the
        chapter, notebook, and task entry point are.
  \item $\Env(s_0)$ is high because posture and location both stabilize rest.
  \item $U(s_1)$ is high if the study task is not pre-specified.
\end{enumerate}
Packing the desk, opening the book in advance, and writing a first micro-task
all raise $P$ and shrink $\DV$ before the transition even begins.
\end{example}

\section{Notation as a Management Tool}

Notation is not cosmetic. It serves two management functions:
\begin{enumerate}[nosep]
  \item it turns vague complaints into diagnosable variables;
  \item it lets multiple writers and reviewers talk about the same failure
        pattern without renaming the problem each time.
\end{enumerate}

\section{Recommended Reading}

\begin{readingbox}
\textbf{Foundation}
\begin{itemize}[nosep]
  \item Eugene Izhikevich, \emph{Dynamical Systems in Neuroscience}.
\end{itemize}
\textbf{Modern Research}
\begin{itemize}[nosep]
  \item The attractor and integrator review literature in computational
        neuroscience provides the clearest bridge from notation to mechanism.
\end{itemize}
\textbf{Frontier}
\begin{itemize}[nosep]
  \item Operator-based or free-energy inspired formalisms should be treated as
        extensions, not prerequisites, for this chapter.
\end{itemize}
\end{readingbox}

\begin{summarybox}
The first chapter turns the manuscript from rhetoric into a controllable
language. Once state, inertia, and barrier are named, later chapters can ask
sharper questions: where is the barrier high, when does it drop, and which
inputs change it most cheaply?
\end{summarybox}
